% 00_abstract.tex  --  Abstract body only (no \begin{abstract} wrapper; main.tex provides that)
% Source: PAPER_SPEC.md §0 locked abstract, light polish only.

Data condensation replaces a large training set with a small surrogate, but
tabular condensers are usually judged only after retraining a model on each
candidate. For histogram gradient-boosted decision trees (GBDTs), we ask whether
a candidate preserves the split decisions the full-data booster would make. We
introduce Split-Landscape Discrepancy (SLD), a no-candidate-retraining
diagnostic that compares full and condensed split-gain landscapes under fixed
preprocessing and prediction state.

In a 17-method audit on 26 datasets using two budgets and three seeds,
method-level split-regret is strongly rank-correlated with downstream AUROC in
the small-to-mid-size regime (Spearman $\rho=-0.87$); a separate five-method,
five-budget, five-seed core grid supports the headline performance comparisons.
SLD is therefore best used as a rejection filter and method-family diagnostic,
not as a single-dataset top-1 selector. It also exposes why coverage can fail:
a greedy high-gain threshold-coverage selector is among the weakest performers
because unweighted coverage differs from gain-weighted structural fidelity. At
very large scale, random sampling becomes hard to beat and root-regret loses
resolution as reasonable surrogates clear dominant split margins. We release
the benchmark, scripts, processed tables, and diagnostic.
